\section{Parameters and Integration}
\label{sec:parameters}

This section explains how parameters interact across the ARC-SPRUCE stack: the Ajtai commitment, the BB-tran rational
hash, the lattice preimage sampler, SmallWood, and the PRF.  Detailed parameter tables and extended benchmark data are
deferred to Appendix~\ref{app:ext-params}.

\subsection{Parameter dependencies}
All components share the same base ring $\Rq$ and modulus $q$:
\begin{itemize}[leftmargin=1.25em]
  \item the commitment matrix, signature map $\mA$, and hash key $B$ are defined over $\Rq$;
  \item SmallWood constraints are evaluated over the base field $\Fq$ and therefore use the same modulus;
  \item the PRF is instantiated over $\Fq$ so that its constraints are native to the proof system.
\end{itemize}
The coefficient bound $\BoundMsg$ governs the message space, holder-side randomness, issuer-side randomness,
centering map, and membership polynomials in the NIZKs.

\subsection{Commitment and shared-randomness parameters}
The Ajtai commitment binds the vector
\[
  [\vecmu\Vert r_0^H\Vert r_1^H\Vert\bar r]^\top.
\]
Its hiding parameters must satisfy the leftover-hash condition of Lemma~\ref{lem:ajtai-hiding}.  The
commitment-only randomness $\bar r$ and any additional committed randomness coordinates must provide enough
min-entropy to statistically hide the full committed vector.

The centering map uses the interval $[-\BoundMsg,\BoundMsg]$ and modulus
\[
  \Delta=2\BoundMsg+1.
\]
For correctness of the shared-randomness step, both holder and issuer contributions must lie in
$[-\BoundMsg,\BoundMsg]$ coefficient-wise.  Lemma~\ref{lem:center-uniform} is then applicable coefficient-wise:
for fixed holder randomness, the centered sum is uniform when the issuer contribution is uniform.

A second entropy budget is needed for target hiding.  Write
\[
  r_0\in\Rq^{\ell_{r_0}}
\]
for the long BB-tran target-randomizer.  Let
\[
  X^N+1\equiv \prod_{j=1}^{f_{\mathrm{fac}}}\Phi_j(X)\pmod q,
  \qquad
  \deg\Phi_j=d_{\mathrm{fac}},
  \qquad
  N=f_{\mathrm{fac}}d_{\mathrm{fac}}.
\]
If the coefficients of $r_0$ are uniform over $[-\BoundMsg,\BoundMsg]$, the bounded-uniform cyclotomic
leftover-hash bound gives
\[
  \varepsilon_{r_0}
  =
  \frac12
  \sqrt{
    \left(
      \left(
        \frac{q^n}{(2\BoundMsg+1)^{\ell_{r_0}}}
      \right)^{d_{\mathrm{fac}}}
      +1
    \right)^{f_{\mathrm{fac}}}
    -1
  }.
\]
The parameters must be chosen so that $\varepsilon_{r_0}\le 2^{-\lambda}$.  Under this condition,
$(B_2,B_2r_0)$ is statistically close to $(B_2,U)$ with $U\getsunif\Rq^n$, and therefore
\[
  T=B_0+B_1\vecmu+B_2r_0+(B_3-\mathbf 1_n r_1)^{-1}
\]
statistically hides $\vecmu$ for every fixed admissible $r_1$; see Lemma~\ref{lem:htran-lhl} and
Theorem~\ref{thm:htran-target-hiding}.  The inverse-side term is treated only as an additive offset, so the
hiding proof uses no bounded-noise assumption on $(B_3-\mathbf 1_n r_1)^{-1}$.

\subsection{Sampler parameters}
The trapdoor sampler must output short signatures $\vecu$ with overwhelming probability and with a distribution close
to a discrete Gaussian over the appropriate lattice coset.  This requires choosing:
\begin{itemize}[leftmargin=1.25em]
  \item ring degree $N$ and modulus $q$,
  \item trapdoor quality,
  \item Gaussian widths for the sampler, and
  \item an acceptance predicate that yields the public verification bound $\BoundSig$.
\end{itemize}
We use an NTRU-style trapdoor together with a Gaussian preimage sampler; details and tuning rules are given in
Appendix~\ref{app:gaussian-sampler}.

\subsection{BB-tran rational-hash parameters}
The BB-tran target is evaluated only when
\[
  B_3-\mathbf 1_n r_1\in\Rq^{\times n}.
\]
Honest target derivation therefore restarts when this condition fails.  The restart probability must be negligible or
or explicitly accounted for as the rational-hash admissibility error $\delta$.  The two BB-tran randomness blocks
play different roles: the long linear-side randomizer $r_0\in\Rq^{\ell_{r_0}}$ carries the entropy required by
the cyclotomic-LHL hiding theorem, whereas the inverse-side randomizer $r_1$ is controlled only by
admissibility and restart probability.  No entropy lower bound on $r_1$ is used in the target-hiding proof.

The pre-sign proof certifies the witness form
\[
  (B_3-\mathbf 1_n r_1)\Had Z=1,
  \qquad
  T=B_0+B_1\vecmu+B_2r_0+Z.
\]
The showing proof certifies the same relation with $T=\mA\vecu$.  Equivalently, it certifies
\[
  (B_3-\mathbf 1_n r_1)\Had\bigl(\mA\vecu-B_0-B_1\vecmu-B_2r_0\bigr)=1.
\]
Membership checks for $[-\BoundMsg,\BoundMsg]$ are implemented via vanishing polynomials over $\Fq$, so one
requires $\BoundMsg\ll q$ to avoid modular wrap-around artifacts.

\paragraph{Non-blind reduction parameters.}
The non-blind reduction imported from DKLW25 is parameterized by a norm bound $\beta_{\mathrm{DKLW}}$, an
inadmissibility bound $\delta$, and the predicate/linear-independence side conditions of the BB-tran family.  The
public verifier bound $\BoundSig$ should therefore dominate the relevant DKLW norm bound under the chosen norm
comparison.

\subsection{SmallWood parameters}
SmallWood introduces parameters controlling packing, masking, and soundness:
\begin{itemize}[leftmargin=1.25em]
  \item the packing size $s=|\Omega|$,
  \item the tail length $\ell$ for HVZK masking and degree enforcement,
  \item the spot-check count $\ell'$ and the number of masks $\rho$ for PIOP soundness,
  \item the degree bounds for witness rows and constraint-derived masks, and
  \item the small-field knob $\theta$ when using an extension-field variant.
\end{itemize}
The pre-sign proof is dominated by the Ajtai opening, centering equations, inverse-witness check, and BB-tran target
identity.  The showing proof adds the signature shortness gadget, the BB-tran signature equation with
$T=\mA\vecu$, and the PRF checkpoint constraints.  Increasing $\ell_{r_0}$ enlarges the witness width and the
number of linear constraints, but it does not increase the nonlinear degree of the BB-tran relation itself.

\subsection{PRF parameters}
The PRF must be secure at the target security level, efficiently representable in constraints, and compatible with
$\Fq$.  We use a Poseidon2-style permutation with feed-forward and truncation.  Its parameters include the state
width $t$, round counts $(R_F,R_P)$, MDS matrices, round constants, S-box exponent $d$, truncation length
$\ell_{\mathsf{tag}}$, and key space $\mathcal K_{\mathsf{key}}$.  We require
$\ell_{\mathsf{tag}}\ge\lceil\lambda/\log_2(q)\rceil$ as an output-length condition; pseudorandomness of the
resulting tags is exactly Assumption~\ref{ass:poseidon-prf} for keys sampled from
$\mathcal K_{\mathsf{key}}$.  The output-length condition alone is not a proof of PRF security.

\subsection{Sizing discussion}
Concrete proof sizes depend jointly on the sampler parameters, the SmallWood masking parameters, the BB-tran relation,
and the PRF checkpoint schedule.  The issuance proof is smaller because the target $T$ is public and the statement only
checks consistency of the committed witness with that target.  The showing proof additionally proves knowledge of a
short preimage and a correct PRF evaluation.
